\RequirePackage[T1]{fontenc}
\documentclass[10pt,conference]{IEEEtran}
\usepackage{amsmath,amssymb,amsthm,graphicx,booktabs,array,url,flushend}
\usepackage[hidelinks]{hyperref}
\usepackage{tikz}
\usetikzlibrary{arrows.meta,positioning}
\emergencystretch=2em
\newtheorem{theorem}{Theorem}
\newtheorem{lemma}{Lemma}
\newtheorem{proposition}{Proposition}
\theoremstyle{definition}
\newtheorem{definition}{Definition}
\newtheorem{remarknum}{Remark}
\newcommand{\code}[1]{\texttt{#1}}
\newcommand{\G}[1]{G#1}
\newcommand{\tool}{ReleaseWeave}
\newcommand{\lang}{RIL}
\begin{document}
\title{Recovering Versioned Releases Across Build and Publish Boundaries}
\maketitle
\begin{abstract}
Distributed release execution must carry dependency outputs as well as source
code. An npm consumer needs its provider's build outputs; an isolated
registry-based Docker builder needs its base image published first. Our
protocol combines stage-aware scheduling with consumer-owned release
obligations, preserving unfinished work after a provider succeeds.
Release planning is a pure function of fixed history, records, graph, and
configuration. A proposed execution profile adds workers to single
repositories, monorepos, and polyrepos. Nodes share one configuration schema
and default to orchestrator capability; explicit workers cannot initiate
releases. The CI-entry node locks all participating repositories and assigns
planned tasks to linked workers within concurrency limits. Temporary Git
branches synchronize changes; verified output transfer lets remote JavaScript
consumers reuse provider builds without rebuilding their prerequisite
closures. Publication retains dependency gates and source-owned records.
The single-history model gives eight conditional recovery guarantees.
All 120 controlled fault episodes recover after repair; an analytical
source-discharge alternative loses unfinished work in 84 (70\%). This compares
admission rules, not competing tools or production failure rates. These results
support the evaluated recovery mechanism; distributed execution requires
implementation and separate performance validation. Full implementation
refinement and external exactly-once effects remain outside the claims.
\end{abstract}
\begin{IEEEkeywords}
software release engineering, continuous integration and delivery,
dependency scheduling, recovery, idempotence
\end{IEEEkeywords}
\input{introduction}
\input{system}
\input{model}
\input{guarantees}
\input{evaluation}
\input{discussion}
\clearpage
\IEEEtriggeratref{34}
\input{references}
\end{document}
